% Appendix A, Study 1 in detail: version 1 and the audit, the achieved panel, the 18 registered comparisons,
% the post-hoc accounting and the report's limitations. The tables are the Study 1 report's, unchanged in
% every entry; the report (main.tex) stands on its own as a posted document. Input from paper.tex; every
% number traces to sources.md and is asserted by check_numbers.py.
\providecommand{\lm}{\texttt{local\_mean\_28}}
\providecommand{\aets}{\texttt{autoets\_weekly}}
\providecommand{\qs}{QScore9}

\section{Study 1 in detail}
\label{app:study1}
\label{sec:results}

\paragraph{Version 1 and the audit.} A first version of the study, built in June and July 2026 on 208
repositories found by three GitHub Search API queries, found two foundation models ahead of four classical
baselines on mean CRPS in every cell of its design. An audit on 2026-07-12 found no code defect that
invalidated the aggregate comparison and two design flaws that did: the cohorts had been selected on
outcomes not available at forecast time, and the evaluation windows mostly tested post-peak decay, the global
peak preceding the evaluated span for 78.0 percent of ``breakout'' repositories. After the audit added a
28-day local mean on identical windows, the foundation models still had the lower mean in every cell while
their per-repository win rate against it fell to between 40.6 and 80.5 percent. The audit's prescription,
freeze and pre-register, context-only labels, trailing non-overlapping windows with a common endpoint, strong
pre-registered baselines, native quantile metrics and repository-clustered inference, became the registered
design of Study 1. Every version 1 number is quoted from the preserved version 1 documents, not recomputed
(Appendix~\ref{app:v1}, Table~\ref{tab:v1}).

\paragraph{The achieved panel.} The frame held 148,550 eligible repositories: 100,687 in the 20 to 99
stratum, 46,811 in 100 to 999, and 1,052 in the 1,000-plus stratum (Table~\ref{tab:panel}). The seeded walk
accepted 100 per stratum and skipped 470 repositories, all because the stable ID no longer resolved to a
public repository at fetch time (466 with HTTP 404) or was forbidden (4 with HTTP 403); the panel is
therefore repositories still public in September 2026 with complete retrievable histories. The daily signal
holds 289,316 repository-days with every Los Angeles calendar day from creation through 2026-08-31 observed;
the creation dates run from 2023-01-02 to 2025-02-19, so every repository was created at least nine months
before the first cutoff, and nothing here is a launch forecast. Over the 900 repository-cutoff windows the
context-only regime labeler assigned 320 windows from 140 repositories to ``inactive recently'', 572 windows
from 228 repositories to ``other active'', and 8 windows from 7 repositories to ``recent surge''; the surge
regime is too thin to say anything about bursts.

\begin{table}[htbp]
\centering\footnotesize
\caption{The achieved panel by day-90 stratum: the frame size, the accepted and skipped counts, the median
and the 10th and 90th percentiles of the survivor ratio, the number of repositories with no star from
2026-06-03 to 2026-08-31, and the median star count over those 90 days. The survivor ratio divides the
endpoint's first 90 Los Angeles days by the gross archive count over the UTC creation day and the next 89
days; the two calendars differ, so the ratio is a diagnostic, not an exact retention rate. Skips are
repositories encountered before the stratum's stop.}
\label{tab:panel}
\begin{tabular}{lrrrrrrrr}
\toprule
Stratum & Frame & Accepted & Skipped & Ratio median & P10 & P90 & Zero, last 90 & Median, last 90 \\
\midrule
20--99 & 100,687 & 100 & 98 & 0.968 & 0.723 & 1.075 & 52 & 0 \\
100--999 & 46,811 & 100 & 348 & 1.004 & 0.885 & 1.079 & 22 & 4 \\
1000+ & 1,052 & 100 & 24 & 1.006 & 0.914 & 1.049 & 1 & 39 \\
\bottomrule
\end{tabular}
\end{table}

\paragraph{The 18 registered comparisons.} Tables~\ref{tab:lm} and~\ref{tab:ae} give the family, one table
per reference, with the mean difference and its interval, the median difference, the number of repositories
where the model's \qs{} is lower, the strict win fraction and the adjusted $p$-value side by side;
Figure~\ref{fig:paired} draws them. Ratio skill against the local mean is defined for only 209, 239 and 259 of
300 repositories at 7, 30 and 90 days, because the reference's loss is zero on the others, so no ratio is
headlined. Pooled \qs{} by stratum (Table~\ref{tab:stratum}) holds stratum means only; a post-hoc reading of
the per-repository accounting file gives the typical repository by stratum: the median difference against
the local mean favors the reference in every stratum and every cell, and the models score lower on 5 to 18 of
100 repositories in the smallest stratum, 10 to 34 in the middle one and 26 to 38 in the top one. Inclusive
80 percent coverage of the repaired deciles over the panel is 0.824 to 0.837 for Chronos-2, 0.880 to 0.910
for TimesFM 3, 0.906 to 0.924 for TimesFM 2.5, 0.941 to 0.955 for the local mean and 0.986 to 0.988 for
AutoETS (Table~\ref{tab:calibration}); on discrete counts a correct quantile can overcover \citep{czado2009},
so these are descriptive. AutoETS's median 90-day daily interval width is 0.79 stars per day in the smallest
stratum, 3.60 in the middle one and 49.69 in the top one (Table~\ref{tab:autoets}).

\begin{table}[htbp]
\centering\small
\caption{The nine declared comparisons against \lm{}. Differences are model minus reference \qs{}, averaged
over three windows within repository and then over 300 repositories; negative favors the model. The interval
is the pointwise 95 percent percentile interval from 20,000 repository-cluster bootstrap replicates. ``Lower''
counts repositories (of 300) where the model's score is strictly lower; the strict win fraction divides by
300. $p$ is the Benjamini-Hochberg adjusted two-sided Wilcoxon signed-rank $p$-value over the 18 tests.}
\label{tab:lm}
\begin{tabular}{llrcrrrr}
\toprule
Model & $h$ & Mean & 95\% interval & Median & Lower & Win fraction & Adjusted $p$ \\
\midrule
Chronos-2 & 7 & -0.137 & [-0.342, 0.009] & 0.0010 & 79 & 0.263 & $1.0 \times 10^{-5}$ \\
Chronos-2 & 30 & -0.132 & [-0.285, -0.022] & 0.0018 & 59 & 0.197 & $5.5 \times 10^{-12}$ \\
Chronos-2 & 90 & -0.101 & [-0.211, -0.016] & 0.0023 & 50 & 0.167 & $1.4 \times 10^{-13}$ \\
TimesFM 3 & 7 & -0.258 & [-0.718, 0.012] & 0.0003 & 89 & 0.297 & $6.7 \times 10^{-3}$ \\
TimesFM 3 & 30 & -0.201 & [-0.488, -0.024] & 0.0004 & 75 & 0.250 & $8.4 \times 10^{-7}$ \\
TimesFM 3 & 90 & -0.128 & [-0.296, -0.016] & 0.0032 & 54 & 0.180 & $6.4 \times 10^{-12}$ \\
TimesFM 2.5 & 7 & -0.243 & [-0.695, 0.025] & 0.0015 & 68 & 0.227 & $3.2 \times 10^{-8}$ \\
TimesFM 2.5 & 30 & -0.201 & [-0.510, -0.015] & 0.0028 & 49 & 0.163 & $2.4 \times 10^{-13}$ \\
TimesFM 2.5 & 90 & -0.112 & [-0.285, 0.001] & 0.0041 & 46 & 0.153 & $1.9 \times 10^{-14}$ \\
\bottomrule
\end{tabular}
\end{table}

\begin{table}[htbp]
\centering\small
\caption{The nine declared comparisons against \aets{}, with the same columns as Table~\ref{tab:lm}.}
\label{tab:ae}
\begin{tabular}{llrcrrrr}
\toprule
Model & $h$ & Mean & 95\% interval & Median & Lower & Win fraction & Adjusted $p$ \\
\midrule
Chronos-2 & 7 & -1.069 & [-1.459, -0.700] & -0.1658 & 296 & 0.987 & $2.1 \times 10^{-48}$ \\
Chronos-2 & 30 & -2.127 & [-2.831, -1.536] & -0.2128 & 298 & 0.993 & $3.3 \times 10^{-49}$ \\
Chronos-2 & 90 & -6.724 & [-13.424, -2.868] & -0.2827 & 298 & 0.993 & $3.1 \times 10^{-49}$ \\
TimesFM 3 & 7 & -1.191 & [-1.547, -0.890] & -0.1817 & 297 & 0.990 & $2.2 \times 10^{-49}$ \\
TimesFM 3 & 30 & -2.196 & [-2.915, -1.583] & -0.2171 & 298 & 0.993 & $8.8 \times 10^{-50}$ \\
TimesFM 3 & 90 & -6.751 & [-13.476, -2.871] & -0.2844 & 293 & 0.977 & $1.9 \times 10^{-48}$ \\
TimesFM 2.5 & 7 & -1.175 & [-1.534, -0.874] & -0.1755 & 297 & 0.990 & $2.2 \times 10^{-49}$ \\
TimesFM 2.5 & 30 & -2.195 & [-2.932, -1.580] & -0.2121 & 298 & 0.993 & $8.8 \times 10^{-50}$ \\
TimesFM 2.5 & 90 & -6.735 & [-13.339, -2.871] & -0.2809 & 297 & 0.990 & $3.2 \times 10^{-49}$ \\
\bottomrule
\end{tabular}
\end{table}

\begin{figure}[htbp]
\centering
\includegraphics[width=\textwidth]{fig-paired-differences}
\caption{The 18 declared comparisons. Filled circles: mean difference (model minus reference \qs{}) with its
pointwise 95 percent repository-bootstrap interval. Open diamonds: median difference. The number under each
interval is the count of repositories, of 300, where the model's score is lower. Note the different vertical
scales.}
\label{fig:paired}
\end{figure}

\paragraph{The post-hoc accounting.} Computed on 2026-09-14 from the saved forecasts after the registered
results were seen, adding no test: each window's \qs{} is split exactly into three parts, each divided by the
full horizon, the loss on zero-target days, which for nonnegative quantiles is
$\frac{2}{9h}\sum_{y_t = 0}\sum_\tau (1-\tau) q_{t,\tau}$, the loss on positive days from quantiles above the
target (overprediction) and from quantiles below it (underprediction). The parts sum to \qs{} on all 48,600
rows to floating-point tolerance and reproduce the 18 registered mean differences (Table~\ref{tab:posthoc}).
The underprediction part favors the models at 7 days, where for both TimesFM models it is the larger of the
two positive-day parts, and favors the local mean at 90 days. Against AutoETS the gain accrues across most
of the panel and its largest part is overprediction on positive days, with a zero-day part between $-0.768$
and $-0.777$. Figure~\ref{fig:zerofraction} places each repository by its context zero fraction: against the
local mean the large model gains are at low zero fractions and the small model losses cluster at high ones.

\begin{table}[htbp]
\centering\small
\caption{Post-hoc. The mean difference against \lm{} and its three parts, model minus reference; negative
favors the model. The parts sum to the total.}
\label{tab:posthoc}
\begin{tabular}{llrrrr}
\toprule
Model & $h$ & Total & Zero days & Positive, overprediction & Positive, underprediction \\
\midrule
Chronos-2 & 7 & -0.137 & 0.005 & -0.095 & -0.048 \\
Chronos-2 & 30 & -0.132 & 0.005 & -0.156 & 0.019 \\
Chronos-2 & 90 & -0.101 & 0.004 & -0.184 & 0.078 \\
TimesFM 3 & 7 & -0.258 & 0.004 & -0.091 & -0.171 \\
TimesFM 3 & 30 & -0.201 & 0.006 & -0.150 & -0.057 \\
TimesFM 3 & 90 & -0.128 & 0.011 & -0.187 & 0.047 \\
TimesFM 2.5 & 7 & -0.243 & 0.012 & -0.088 & -0.167 \\
TimesFM 2.5 & 30 & -0.201 & 0.014 & -0.163 & -0.052 \\
TimesFM 2.5 & 90 & -0.112 & 0.015 & -0.209 & 0.082 \\
\bottomrule
\end{tabular}
\end{table}

\paragraph{Study 1's limitations, as the report gave them.}
\begin{itemize}
\item \emph{A fetch-time survivor panel.} The 470 skips were repositories that no longer resolved publicly in
  September 2026. Equal weighting describes the panel of repositories still public with complete histories,
  not all repositories created in the window.
\item \emph{Surviving stars, not gross additions.} The endpoint reports stars that still exist at retrieval.
  In the 2024 validation the endpoint total was below the archive total for 17 of 20 repositories; for the
  15 large projects the ratio ran from 0.728 to 0.994, with 13 of 15 at least 5 percent below. The study
  target is the surviving-star history; no correction is applied, and later unstars can change the apparent
  past.
\item \emph{Los Angeles day bins.} The endpoint's day boundary was resolved to America/Los\_Angeles days by
  validation, not documented by GitHub; the validation also showed that the transition days around daylight
  saving do not uniformly favor one calendar.
\item \emph{Three windows per repository, and horizons sharing targets.} Inference is conditional on the
  fixed calendar of three cutoffs; the 7- and 30-day targets are prefixes of the 90-day targets, so the
  nine cells per reference are not independent.
\item \emph{Pointwise intervals and BH.} The bootstrap intervals are pointwise, not simultaneous, and the
  Benjamini-Hochberg adjustment carries its usual dependence assumptions.
\item \emph{The signed-rank symmetry assumption.} The Wilcoxon test's location interpretation requires
  symmetric differences, which these are not; it is reported as registered and is not a test of the mean.
\item \emph{AutoETS as a weak uncertainty reference.} Its 80 percent intervals cover about 99 percent of
  target days; a model can beat it on nearly every repository by being less wide. Read with the coverage
  results, the comparison against the local mean is the more informative of the two; that is a reading made
  after the results, and both comparisons stay reported as registered.
\item \emph{Training contamination unresolved.} No author-disclosed complete-corpus training cutoff exists
  for any of the three models; the model cards name specific sources with dates, which do not bound the
  whole corpus. What can be said is bound by availability, per model. The targets, 2025-12-05 to 2026-08-31,
  postdate the pinned TimesFM 2.5 revision, dated 2025-10-02, but are not shown to postdate the pinned
  Chronos-2 revision, dated 2026-06-05, or the pinned TimesFM 3 revision, dated 2026-09-02 with weights
  uploaded 2026-08-24; these are repository revision dates, not training cutoffs. Every repository's history
  before the cutoffs may be familiar to all three models, and a later test window on familiar entities does
  not establish generalization to unseen ones \citep{moghadasi2026}. The controls are Study 1's registered
  follow-up cohort and the lapsed October holdout; neither has run (Section~\ref{sec:notrun}).
\item \emph{The mature-phase calendar.} Every repository is forecast between nine and 41 months after
  creation. Nothing here concerns launches or the days after a release.
\item \emph{Version 1 not recomputed.} The version 1 numbers are quoted from preserved documents.

\end{itemize}

\paragraph{Study 1's follow-up cohort.} The protocol registers a separate follow-up on repositories created
on or after 2026-09-13, to run after the write-up with identical eligibility, strata, sample size, model
revisions, baselines, scores, regime rules and relative window calendar; it is an availability-bound
exclusion based on the already frozen checkpoints, not a claimed training cutoff, and it has not run.
